\subsection*{Density and the Archimedean Property}

\begin{definitionbox}{Definition}
\begin{definition}[Two in $\mathbb{Q}$]
\label{def:two-in-q}
Define
\[
2_{\mathbb{Q}}:=1_{\mathbb{Q}}+1_{\mathbb{Q}}.
\]
\end{definition}
\end{definitionbox}

\begin{dependencies}
\begin{itemize}
  \item TODO: determine dependencies for \texttt{def:two-in-q}.
\end{itemize}
\end{dependencies}

\begin{lemmabox}{Lemma}
\begin{lemma}[Two Is Nonzero in $\mathbb{Q}$]
\label{lem:two-nonzero-in-q}
In $\mathbb{Q}$,
\[
2_{\mathbb{Q}}\ne 0_{\mathbb{Q}}.
\]
\end{lemma}
\end{lemmabox}

\begin{dependencies}
\begin{itemize}
  \item \hyperref[cor:one-positive-on-q]{One Is Positive in $\mathbb{Q}$}
  \item \hyperref[thm:trichotomy-on-q]{Trichotomy on $\mathbb{Q}$}
\end{itemize}
\end{dependencies}

\begin{lemmabox}{Lemma}
\begin{lemma}[Midpoint Lies Strictly Between Rationals]
\label{lem:midpoint-between-q}
For all $x,y\in\mathbb{Q}$,
\[
x<y
\Longrightarrow
x< (x+y)\cdot 2_{\mathbb{Q}}^{-1}<y.
\]
\end{lemma}
\end{lemmabox}

\begin{dependencies}
\begin{itemize}
  \item \hyperref[thm:q-is-an-ordered-field]{$\mathbb{Q}$ Is an Ordered Field}
  \item \hyperref[lem:two-nonzero-in-q]{Two Is Nonzero in $\mathbb{Q}$}
\end{itemize}
\end{dependencies}

\begin{theorembox}{Theorem}
\begin{theorem}[Density of $\mathbb{Q}$ in Itself]
\label{thm:density-of-q}
For all $x,y\in\mathbb{Q}$,
\[
x<y
\Longrightarrow
\exists r\in\mathbb{Q}\ (x<r<y).
\]
\end{theorem}
\end{theorembox}

\begin{dependencies}
\begin{itemize}
  \item \hyperref[def:order-on-q]{Order on $\mathbb{Q}$}
  \item \hyperref[lem:midpoint-between-q]{Midpoint Lies Strictly Between Rationals}
\end{itemize}
\end{dependencies}

\begin{theorembox}{Theorem}
\begin{theorem}[Archimedean Property of $\mathbb{Q}$]
\label{thm:archimedean-property-of-q}
For all $x,y\in\mathbb{Q}$,
\[
0_{\mathbb{Q}}<x
\Longrightarrow
\exists n\in\mathbb{N}\ (y<n\cdot x).
\]
\end{theorem}
\end{theorembox}

\begin{dependencies}
\begin{itemize}
  \item \hyperref[thm:well-ordering-principle]{Well-Ordering Principle}
  \item \hyperref[thm:division-algorithm-on-z]{Division Algorithm on $\mathbb{Z}$}
  \item \hyperref[thm:z-totally-ordered-set]{$\mathbb{Z}$ Is a Totally Ordered Set}
\end{itemize}
\end{dependencies}

\subsection*{From Discreteness to Arbitrary Closeness}

\begin{remark}[Discrete Systems Before $\mathbb{Q}$]
\label{rem:discrete-systems-before-q}
The systems $\mathbb{N}$, $\mathbb{W}$, and $\mathbb{Z}$ are discrete in the
order sense: consecutive integer points have no same-system point between
them. The rational numbers are the first number system in this construction
chain where intervals can be subdivided internally without leaving the system.
\end{remark}

\begin{definitionbox}{Definition}
\begin{definition}[Arbitrarily Close Distinct Points]
\label{def:arbitrarily-close-distinct-points}
A number system has arbitrarily close distinct points if for every point $x$
and every positive tolerance $\varepsilon$, there is a point $y\ne x$ such that
\[
|x-y|<\varepsilon.
\]
\end{definition}
\end{definitionbox}

\begin{dependencies}
\begin{itemize}
  \item TODO: determine dependencies for \texttt{def:arbitrarily-close-distinct-points}.
\end{itemize}
\end{dependencies}

\begin{theorembox}{Theorem}
\begin{theorem}[$\mathbb{Q}$ Has Arbitrarily Close Distinct Points]
\label{thm:q-has-arbitrarily-close-distinct-points}
For every $x\in\mathbb{Q}$ and every $\varepsilon\in\mathbb{Q}$,
\[
0_{\mathbb{Q}}<\varepsilon
\Longrightarrow
\exists y\in\mathbb{Q}\ (y\ne x\ \text{and}\ |x-y|<\varepsilon).
\]
\end{theorem}
\end{theorembox}

\begin{dependencies}
\begin{itemize}
  \item \hyperref[cor:reciprocals-arbitrarily-small-q]{Reciprocals Become Arbitrarily Small in $\mathbb{Q}$}
  \item \hyperref[thm:addition-preserves-strict-order-on-q]{Addition Preserves Strict Order on $\mathbb{Q}$}
\end{itemize}
\end{dependencies}

\begin{corollarybox}{Corollary}
\begin{corollary}[$\mathbb{Q}$ Has No Adjacent Points]
\label{cor:q-has-no-adjacent-points}
There do not exist rational numbers $x<y$ with no rational number strictly
between them.
\end{corollary}
\end{corollarybox}

\begin{dependencies}
\begin{itemize}
  \item \hyperref[thm:density-of-q]{Density of $\mathbb{Q}$ in Itself}
\end{itemize}
\end{dependencies}

\begin{propositionbox}{Proposition}
\begin{proposition}[$\mathbb{Q}$ Has No Least Positive Rational]
\label{prop:no-least-positive-rational}
For every $x\in\mathbb{Q}$,
\[
0_{\mathbb{Q}}<x
\Longrightarrow
\exists y\in\mathbb{Q}\ (0_{\mathbb{Q}}<y<x).
\]
\end{proposition}
\end{propositionbox}

\begin{dependencies}
\begin{itemize}
  \item \hyperref[thm:density-of-q]{Density of $\mathbb{Q}$ in Itself}
\end{itemize}
\end{dependencies}

\begin{corollarybox}{Corollary}
\begin{corollary}[Reciprocals Become Arbitrarily Small in $\mathbb{Q}$]
\label{cor:reciprocals-arbitrarily-small-q}
For every $\varepsilon\in\mathbb{Q}$,
\[
0_{\mathbb{Q}}<\varepsilon
\Longrightarrow
\exists n\in\mathbb{N}\ 
\left(0_{\mathbb{Q}}<n^{-1}<\varepsilon\right).
\]
\end{corollary}
\end{corollarybox}

\begin{dependencies}
\begin{itemize}
  \item \hyperref[thm:archimedean-property-of-q]{Archimedean Property of $\mathbb{Q}$}
  \item \hyperref[thm:sign-of-reciprocal-q]{Sign of a Reciprocal on $\mathbb{Q}$}
\end{itemize}
\end{dependencies}

\subsection*{Limit-Like Behavior Before Limits}

\begin{propositionbox}{Proposition}
\begin{proposition}[One-Sided Rational Approximation]
\label{prop:one-sided-rational-approximation}
For every $x\in\mathbb{Q}$ and every $\varepsilon\in\mathbb{Q}$,
\[
0_{\mathbb{Q}}<\varepsilon
\Longrightarrow
\exists y,z\in\mathbb{Q}\ (x-\varepsilon<y<x<z<x+\varepsilon).
\]
\end{proposition}
\end{propositionbox}

\begin{dependencies}
\begin{itemize}
  \item \hyperref[thm:q-has-arbitrarily-close-distinct-points]{$\mathbb{Q}$ Has Arbitrarily Close Distinct Points}
  \item \hyperref[thm:density-of-q]{Density of $\mathbb{Q}$ in Itself}
\end{itemize}
\end{dependencies}

\begin{remark}[Refinement and Bounds]
\label{rem:rational-refinement-and-bounds}
The Archimedean property supplies arbitrarily small rational increments, so a
rational point can be approached from one side at any prescribed tolerance.
Mediants supply a complementary order-refinement mechanism: once two rational
bounds are known, their mediant is a new rational bound strictly between them.
The following approximation topics make these two mechanisms explicit through
dyadic grids, mediant refinement, Stern--Brocot paths, and rational Cauchy
sequences.
\end{remark}
